% STATUS: verified
\chapter{Introduction}\label{ch:intro}

This chapter says what the \product{} is for, what it does, and where in it a language model
is at work. It ends with the system in numbers, the course it was built for, and what the
system does not claim to do.
Everything it mentions is explained in a later chapter; the references say where.

\section{The problem}\label{sec:intro:problem}

A financial institution in Lithuania answers to several dozen legal acts at once: regulations
of the European Union that bind it directly, directives that reach it through national law,
Lithuanian laws, and the resolutions, instructions and guidelines of its supervisors. A
crypto-asset service provider, for instance, must find its duties in the regulation on markets
in crypto-assets and its technical standards, in the regulation on digital operational
resilience, in the rules on money-laundering prevention and on transfers of funds, in the
data-protection regulation, and in the national laws that give all of these their fines. The
person responsible for compliance has to establish four things, and to establish them again
whenever a law or the company changes:

\begin{enumerate}
  \item \textbf{Which acts apply} to this company, given what it is and what it does.
  \item \textbf{What each of them obliges} the company to do, provision by provision.
  \item \textbf{Whether the company's own documents} --- its policies and procedures --- carry
        each obligation out, and where they fall short.
  \item \textbf{What to do about the gaps}, in what order, and who owns each action.
\end{enumerate}

Done by hand this is weeks of reading. The reading is also exactly the kind of work a language
model does quickly and, left alone, unreliably: it will cite an article that says something
else, quote a sentence the policy does not contain, or call a duty met because the policy
mentions the subject. The \product{} is built around that tension. It lets models do the
reading, and it surrounds their calls with three things: the text of the law itself,
fetched from its official source, so that the law is read and not recalled from memory; checks made by
ordinary code, so that a quotation or a citation is verified and not trusted; and a person who
decides, wherever a decision has consequences.

\section{What the system does}\label{sec:intro:does}

The \product{} (\lt{\productlt} in its first language, Lithuanian) is a self-hosted web
application. A user signs in and lands in a conversation with the \term{chat assistant}; pages
exist beside the chat for reviewing results and for administration. In the order in which a
user meets them, the system:

\begin{description}
  \item[Keeps the law as text.] A \term{legal corpus} of 78 catalogued acts, 68 of them in
    use, is fetched from the official machine sources of the European Union and of Lithuania
    (two guidelines are uploaded as files), parsed down to article, paragraph and point, and
    kept in versions. Findings and answers cite provisions of this corpus and can show their
    words (\cref{ch:corpus}).
  \item[Turns the law into a register of obligations.] Each act is read once and the duties it
    imposes are written down as \term{obligations} --- who owes what, under which condition, by
    when, how gravely --- each tied to the provisions it comes from. A second model call then checks
    every obligation against its provision, and the design leaves the approval of each to a
    person (\cref{ch:obligations}; how the register that ships was reviewed is stated in
    \cref{ch:control}).
  \item[Gets to know a company.] The user describes the company in the chat and drops in its
    policies; the assistant files the profile and the documents through its tools. Each
    document is indexed for search and read into a structured \term{inventory} of what it
    establishes (\cref{ch:company}).
  \item[Answers questions about the law.] The chat searches the corpus, quotes the exact
    provision, and checks whether an article a document cites really regulates what the
    document says it does (\cref{ch:chat,ch:company}).
  \item[Makes a compliance report.] An analysis decides which acts apply to the company,
    stops for a person to approve that list (unless the person who started it chose to let it
    run through), then compares every obligation of those acts that fits the company with the company's documents. The result is a \term{report}: for every such obligation a \term{finding}
    with its status, the evidence quoted from the documents, a recommendation and, for a gap
    the act sanctions, what a breach can cost; an \term{action plan} that covers every gap; and a summary for the board.
    A person reviews the findings and signs the report off. It exports to PDF and Excel
    (\cref{ch:report}).
  \item[Follows the actions over time.] The plan of a signed-off report moves into the
    company's \term{action register}, where an action is closed when a later signed-off
    report finds its obligations met, or by a person by hand with a note that says why
    (\cref{ch:register}).
  \item[Drafts the documents of a compliance programme.] From the same documents the system
    drafts three \term{work products}: a compliance policy with the documents that are still
    missing; the governance structure, the compliance process and its responsibility matrix;
    and the legal risk criteria and register. Code checks each against the rules of its kind,
    and a person edits and approves it (\cref{ch:programme}).
  \item[Shows its own workings.] An administrator manages the corpus, reviews obligations,
    reads and edits every prompt, and sees every recorded model call with its cost, and every change
    anyone made (\cref{ch:console,ch:prompts,ch:security}).
\end{description}

\section{From the law to the report in one picture}\label{sec:intro:picture}

\Cref{fig:intro:map} shows the whole system as four tracks. The first runs once per act, long
before any company is involved: it is the costly, careful reading of the law, and every report
reuses it. The second reads each document once; only the check of its citations is repeated
by every report. The third is a report; the fourth a work product. The chat assistant stands
above them: through its tools it reads all four, files a company's documents and starts a
report or a work product, and it never does their work itself.

\begin{figure}[H]
\centering
\begin{tikzpicture}[
    n/.style={text width=15.5mm, inner xsep=3pt, inner ysep=4pt, minimum height=11.5mm, font=\sffamily\scriptsize},
    node distance=5mm and 3.2mm]
  % --- the law
  \node[ext, n] (src) {official sources};
  \node[svc, n, right=of src] (corp) {legal corpus: acts, provisions};
  \node[ai, n, right=of corp] (ext) {obligations extracted};
  \node[ai, n, right=of ext] (jud) {a judge checks each};
  \node[man, n, right=of jud] (rev) {a person reviews};
  \node[store, n, right=of rev, text width=15mm] (reg) {register of obligations};
  \draw[flow] (src) -- (corp); \draw[flow] (corp) -- (ext); \draw[flow] (ext) -- (jud);
  \draw[flow] (jud) -- (rev); \draw[flow] (rev) -- (reg);
  % --- the company
  \node[data, n, below=13mm of src, text width=13mm] (docs) {a company's documents};
  \node[ai, n, right=of docs] (idx) {indexed: embeddings};
  \node[ai, n, right=of idx] (inv) {inventory of each document};
  \node[ai, n, right=of inv] (cit) {citations checked};
  \draw[flow] (docs) -- (idx); \draw[flow] (idx) -- (inv); \draw[flow] (inv) -- (cit);
  % --- the report
  \node[ai, n, below=13mm of docs] (app) {which acts apply};
  \node[man, n, right=of app] (scope) {scope approved};
  \node[ai, n, right=of scope] (gap) {gap analysis: findings};
  \node[chk, n, right=of gap] (chk) {quotes located, sanctions joined};
  \node[ai, n, right=of chk] (plan) {action plan, summary};
  \node[man, n, right=of plan] (sign) {reviewed, signed off};
  \node[chk, n, right=of sign] (areg) {action register};
  \draw[flow] (app) -- (scope); \draw[flow] (scope) -- (gap); \draw[flow] (gap) -- (chk);
  \draw[flow] (chk) -- (plan); \draw[flow] (plan) -- (sign); \draw[flow] (sign) -- (areg);
  % --- work products
  \node[ai, n, below=13mm of app] (wp) {policy, structure, risk register drafted};
  \node[chk, n, right=of wp] (rules) {rules of the kind checked};
  \node[man, n, right=of rules] (appr) {edited, approved};
  \draw[flow] (wp) -- (rules); \draw[flow] (rules) -- (appr);
  % --- joins
  \draw[opt] (reg.south) |- ($(gap.north)+(0,5mm)$) -- (gap.north);
  \draw[opt] (cit.south) -- ++(0,-4mm) -| ($(gap.north)+(-3mm,0)$);
  \draw[opt] (docs.west) -- ++(-3mm,0) |- (wp.west);
  % --- lanes
  \begin{scope}[on background layer]
    \node[lane, fit=(src)(reg)] (l1) {};
    \node[lane, fit=(docs)(cit)] (l2) {};
    \node[lane, fit=(app)(areg)] (l3) {};
    \node[lane, fit=(wp)(appr)] (l4) {};
  \end{scope}
  \node[lanelabel, anchor=south west] at (l1.north west) {The law --- once per act};
  \node[lanelabel, anchor=south west] at (l2.north west) {The company --- its documents};
  \node[lanelabel, anchor=south west] at (l3.north west) {A report};
  \node[lanelabel, anchor=south west] at (l4.north west) {A work product};
  % --- the chat
  \node[ai, n, text width=30mm, right=6mm of appr, minimum height=11.5mm] (chat) {the chat assistant reads all of this, and starts reports and work products, through its tools};
\end{tikzpicture}
\actorlegend
\caption{The system as four tracks. Yellow boxes are calls to a model, blue boxes are
checks or computations by code, dark boxes are decisions of a person. The dotted lines show
what a report takes from the tracks above it: the obligations of the acts that apply, and what
the company's documents say.}\label{fig:intro:map}
\end{figure}

Three rules, visible in the picture, run through the whole system and are the subject of
\cref{ch:control}:

\begin{enumerate}
  \item \textbf{The law is read, not remembered.} A model that reads an act is given the
    provision's text from the corpus and asked a question about it, and a report judges a
    company against the obligations read that way. Where a structured answer names a
    provision, code looks the provision up; where it quotes a document, code finds the
    quotation in the document. The chat's own text is the exception: nothing checks it, and
    the chat is told to mark every number it has not read in the corpus as unverified
    (\cref{ch:chat}).
  \item \textbf{The model drafts, code checks, a person decides.} Scores, bands, the order of a plan and
    the fine attached to a finding are computed, never taken from a model. Nothing is ever
    marked approved by the system itself: an obligation of the register, a finding's override,
    the sign-off of a report and a work product are each approved by a named person, and the
    approval is recorded. The scope of a report is approved the same way; a person may instead
    let a report run through without stopping, and the report then records that nobody
    approved its scope. The chat assistant approves nothing on its own. It can record three
    decisions, always in the user's name: a report's scope, on the user's yes and at the
    price the user accepted; a finding's override, on the user's decision; and a credential
    marked verified, on the user's word that it was checked. It approves no obligation, work
    product or report (\cref{ch:chat}).
  \item \textbf{Every call is on the record.} Each model call that answers is logged with its
    tokens, its cost and, when a prompt produced it, the prompt version (only the embedding of
    a search query, a few tokens, is left out on purpose); each prompt is stored in versions an
    administrator can read, change and roll back; each report records which versions of which
    prompts and acts it was made from.
\end{enumerate}

\section{Where the models are}\label{sec:intro:models}

A language model is called in twenty-three places, each with a stored prompt of its own, and a
second kind of model turns text into vectors for search. \Cref{tab:intro:paths} groups them;
Part~II of this document takes them one by one, and \cref{app:prompts} describes every prompt.
All of them are reached through one service and one gateway (\cref{ch:models}).

\begin{table}[H]
\caption{Every place a model is called, by what it is for}\label{tab:intro:paths}
\begin{tabularx}{\textwidth}{@{}L{31mm}YL{19mm}@{}}
\toprule
\hd{What for} & \hd{What the model is asked} & \hd{Chapter} \\
\midrule
Reading the law, once per act &
  To write down the obligations of an act's provisions (\prompt{corpus.extract_obligations}) and
  the requirements of a standard (\prompt{corpus.extract_requirements}); to say whom an act's own
  scope provisions cover (\prompt{corpus.scope}); to read its sanction provisions into tiers
  (\prompt{corpus.sanctions}); to mark the headings of a standard that is not numbered the usual
  way (\prompt{corpus.standard_outline}); and, as a judge, to check each extracted obligation
  against the provision it cites (\prompt{corpus.verify_obligations}). &
  \cref{ch:corpus,ch:obligations} \\
Reading a company's documents &
  To read a document into an inventory of what it establishes (\prompt{report.inventory}); and, as a
  judge, to say whether an article a document cites regulates what the document claims
  (\prompt{corpus.check_citations}). &
  \cref{ch:company} \\
Making a report &
  To decide whether an act applies to the company (\prompt{report.applicability}); to judge each
  obligation against the documents (\prompt{report.gap_analysis}); to read the company's own
  requirements out of its policies and judge them (\prompt{company.internal_requirements},
  \prompt{report.internal_gap}); to plan the actions and merge the drafts
  (\prompt{report.action_plan}, \prompt{report.plan_merge}); to write the board's summary
  (\prompt{report.summary}). &
  \cref{ch:report} \\
Drafting work products &
  To draft the compliance policy and list the missing documents (\prompt{programme.policy},
  \prompt{programme.documents}); the structure and the responsibility matrix
  (\prompt{programme.structure}, \prompt{programme.raci}); the risk criteria and scenarios
  (\prompt{programme.risk_criteria}, \prompt{programme.risk_scenarios}). &
  \cref{ch:programme} \\
Conversation &
  To hold the chat and choose which tools to call (\prompt{chat.assistant}); to give a thread its
  title (\prompt{chat.title}). &
  \cref{ch:chat} \\
Search &
  An embedding model turns the provisions of the corpus, the obligations, the passages of a
  company's documents and each search query into vectors. &
  \cref{ch:models,ch:corpus} \\
\bottomrule
\end{tabularx}
\end{table}

Just as telling is where no model is called. Fetching and parsing the law (but for the
headings of a standard the rules cannot read), locating a quotation, looking a provision up, joining a finding to its sanction, computing a risk score,
ordering a plan, moving the action register, and every approval are code or people. The
action register in particular --- the part of the system that says whether a gap was really
closed --- contains no model call at all (\cref{ch:register}).

\section{The system in numbers}\label{sec:intro:numbers}

The figures below were read on the deployed system on 13~November 2026; the two that have
changed since, the count of prompt versions and the sanction tiers --- read again on a copy
of the corpus with the new version of their prompt --- are given as they stand since
19~November 2026.
They are given so that the reader has a sense of scale; the chapters say how each was obtained.

\begin{table}[H]
\caption{The system on 13 November 2026}\label{tab:intro:numbers}
\begin{tabularx}{\textwidth}{@{}L{44mm}Y@{}}
\toprule
\hd{What} & \hd{How much} \\
\midrule
Legal corpus & 78 acts in the catalogue, 68 of them enabled; 26\,983 provisions stored as units in
  their current versions; 4\,059 passages embedded for search (\cref{app:acts}) \\
Register of obligations & 5\,960 obligations extracted from 66 of the 68 acts (two state none): 5\,938 approved,
  10 rejected, 12 left for a person; 302 sanction tiers read from 27 acts \\
Prompts & 23, in 91 versions since 20~November 2026 (85 until 16~November); the chat's prompt is in
  its seventeenth version, the report summary's in its twelfth \\
One full report & Made on 13~November 2026, without the scope review, for the course's sample company, a crypto-asset service
  provider with 26 policies: 26 acts assessed (one of them as preparation for a date still
  ahead), 2\,151 findings from the law and 475 from the company's own policies, an action plan of 51 actions covering all 1\,633 open findings;
  272 model calls, 1 hour 38 minutes, \$13.20 (\cref{ch:report}) \\
Code & About 33\,600 lines in the API, 10\,700 in the assistant service and 46\,800 in the
  console, with about 94\,000 lines of tests beside them \\
Tests & 405 in the API, 509 in the assistant service, 5\,180 in the console
  (\cref{ch:testing}); six suites of labelled cases that measure six of the prompts against
  the real model, every label made on 12~November 2026 by the AI coding assistant used to
  build the system and none yet confirmed by a person (\cref{ch:control}) \\
\bottomrule
\end{tabularx}
\end{table}

\section{The course the system was built for}\label{sec:intro:course}

The \product{} is the project of the course \lt{Finansų teisė ir finansų įstaigų veiklos
atitikties valdymas} (financial law and the compliance management of financial institutions)
of the FinTech master's programme at Vilnius University Kaunas Faculty, autumn 2026. The course
teaches a method: classify the institution, build its compliance policy and governance, assess
its legal risk, monitor and report. It gives three assignments on a sample organisation,
\demo{} --- a fictional crypto-asset service provider with 26 policies --- and it describes
the role of the technologist who builds tools for compliance work: prompts bound to schemas,
answers grounded in retrieved sources, kept sets of labelled cases, regression measurement,
audit trails. The system is an attempt to build that method as software, and the sample
organisation is the company on which the reports and the other company-side measurements of
this document were made.
\Cref{ch:course} maps the course's assignments and concepts onto the system part by part.

\section{What the system does not claim}\label{sec:intro:limits}

\begin{itemize}
  \item \textbf{It is not legal advice.} A report is a structured first reading for a
    compliance officer to review: its summary carries a section on its own limitations, and it
    records how many of its findings rest on approved obligations of the register.
  \item \textbf{It judges what is written down.} A finding says whether the company's documents
    establish what an obligation asks for --- the design of its arrangements. Whether the
    company does in practice what its policy says is outside what documents can show; a
    finding that is \status{met} carries a note to verify it in practice where that matters.
  \item \textbf{It reads text.} A scanned document without a text layer gives it nothing to
    read: there is no optical character recognition.
  \item \textbf{Its law is the law of its corpus.} It covers the financial regulation of
    the European Union and of Lithuania that the catalogue lists (\cref{app:acts}), and any
    act an administrator adds. An act outside the corpus does not exist for it, and it says
    so when a document cites one.
  \item \textbf{Its measurements are observations.} The accuracy figures, costs and durations
    in this document were measured on one sample company, or on the shared corpus, on the days
    named. They show how the
    system behaved, not how it must behave elsewhere.
\end{itemize}

The limitations of each part are stated in its chapter, and gathered in \cref{ch:control}.
